\section{The gain in the initial box exponent}
\label{sec:area_law_initial_exponent_gain}

Write $C_{\mathrm{pad}}$ for the padding constant used in
Theorem~\ref{thm:area_law_initial_box}. The cyclic-color recurrence uses
$p=1/(2(2C_{\mathrm{pad}}+1)^2)$ and $M=8(C_{\mathrm{pad}}+1)$.
Its ratio is $\lambda=(1-p/2)M^2$, and its exponent is
$e_0=\log\lambda/\log M-1$.
The following auxiliary estimates quantify the gain for these formalized
recurrence constants. The manuscript instead uses
$p_0=(2C_{\mathrm{pad}}+5)^{-2}$ and
$\lambda=(1-p_0/4)M^2$, with $M\ge16(C_{\mathrm{pad}}+2)$;
the numerical choices here come from the cyclic-color partition used
in the initial recurrence.

\begin{lemma}[Exact initial exponent gain]
    \label{thm:area_law_initial_exponent_gain}
    \lean{TNLean.PEPS.AreaLaw.one_sub_initialBoxExponent_eq}
    \leanok
    \uses{thm:area_law_initial_box}
    For the initial recurrence parameters,
    \begin{align}
        1-e_0=\frac{-\log(1-p/2)}{\log M}.
    \end{align}
    This is the logarithmic form of the exponent $e_0$, following the exponent
    calculation in
    \cite[proof of Proposition~3.3, lines~651--663]{OpenAI2026AreaLaw}.
\end{lemma}
\begin{proof}\leanok
    \lean{TNLean.PEPS.AreaLaw.initialBoxProb_pos_and_le_half,
        TNLean.PEPS.AreaLaw.sixteen_le_initialBoxScale,
        TNLean.PEPS.AreaLaw.one_lt_initialBoxScale}
    Since $0<p\le1/2$ and $M\ge16$,
    $\log\lambda=\log(1-p/2)+2\log M$ and $\log M>0$.
    Substitute this identity into $e_0=\log\lambda/\log M-1$.
\end{proof}

\begin{lemma}[Two-sided initial exponent gain]
    \label{thm:area_law_initial_exponent_gain_bounds}
    \lean{TNLean.PEPS.AreaLaw.initialBoxExponent_gain_bounds}
    \leanok
    \uses{thm:area_law_initial_box}
    For the initial recurrence parameters,
    \begin{align}
        \frac{p}{2\log M}\le1-e_0\le\frac{p}{(2-p)\log M}.
    \end{align}
\end{lemma}
\begin{proof}\leanok
    \uses{thm:area_law_initial_exponent_gain}
    Put $x=1-p/2>0$. The inequalities $\log x\le x-1$ and
    $1-x^{-1}\le\log x$ give
    $p/2\le-\log(1-p/2)\le p/(2-p)$.
    Divide by $\log M>0$.
\end{proof}
